\documentclass[12pt]{article}
\usepackage[utf8]{inputenc}
\usepackage[T1]{fontenc}
\usepackage{hyperref}
\title{Ulch Language Resources for K-12 Curriculum Development}
\author{Compiled by Claude Sonnet 5 (claude-sonnet-5) for Andr\'as Kornai}
\date{}
\begin{document}
\maketitle
\section{Introduction}
Ulch (ISO 639-3 ulc) is an endangered Tungusic language of the Russian Far East, closely related to Nanai and Orok and long treated as a Nanai dialect. Search results report 506 speakers in the 2020 census, about 20\% of an ethnic group of 2,481, with use limited to older generations. Few resources exist, and none found is a full curriculum.
\section{Primary Educational Resources}
Search results report that a school writing system was developed in the late 1980s by L.~I.~Sem (Primer, 1992), and that T.~B.~Matveeva's \emph{Ulch Language, Grade 3} appeared in St Petersburg in 2018. They also mention \emph{First Steps in the Ulch Language} (Khabarovsk, 2006), a collection of tongue twisters, poems and tales. I did not find online copies of these; the descriptions come from summaries of the search results. Overview pages: \url{https://lingsib.iea.ras.ru/en/languages/ulch.shtml} and \url{https://atlaskmns.ru/page/en/lang_ulchi_all.html}.
\section{Linguistic and Reference Materials}
Reported works include Schmidt's 1908 materials (published 1923), Lopatin (1957), and Sunik, \emph{The Ulch Language: Studies and Materials} (1985). A 2022 article by V.~A.~Gorbunova discusses principles for a prospective Ulch--Russian dictionary: \url{https://www.philology.nsc.ru/journals/kis/pdf/CS_2022_2/08.pdf}. Journal records: \url{https://philology.nsc.ru/journals/ykns_en/article.php?id=540} and \url{https://philology.nsc.ru/journals/kis_en/article.php?id=696}. Related Nanai overview: \url{https://lingsib.iea.ras.ru/en/languages/nanai.shtml}.
\section{Religious and Community Resources}
No religious or community-organization resources were found in these searches.
\section{Digital and Audio Resources}
A language profile is at \url{https://docs.verbix.com/EndangeredLanguages/Ulch}. No audio or software resources were found.
\section{Curriculum Development Context}
Ulch is not known to be taught as a full school subject beyond the limited materials above. Any curriculum would need new work with elderly speakers and the existing Nanai-based and Sem orthography materials. Little else was found.
\end{document}